% ===========
% CHƯƠNG 10
% ===========
\OpenGrid

\ParallelChapter
  {Coherence}
  {Tính nhất quán}
  {}

\TriPara
  {
    \EN{
      \begin{flushright}
        \emph{Curriculum, instruction, and assessment form a coherent whole to support reasoning and sense making.}
      \end{flushright}
    }
  }
  {
    \VI{
      \begin{flushright}
        \emph{Chương trình học, giảng dạy và đánh giá tạo thành một chỉnh thể nhất quán nhằm hỗ trợ luận và hiểu.}
      \end{flushright}
    }
  }
  {
    \NOTE{Nguyên tắc nền tảng: luận và hiểu chỉ có thể được nuôi dưỡng bền vững khi chương trình học, giảng dạy và đánh giá vận hành như một chỉnh thể nhất quán. Nếu ba thành tố này không đi cùng hướng, thì nỗ lực thúc đẩy luận và hiểu trong lớp học sẽ dễ bị suy yếu hoặc triệt tiêu.\par\vspace*{\baselineskip}

    \textbf{Ba khái niệm then chốt}\par\vspace*{\baselineskip}

    \begin{itemize}
      \item \emph{Chương trình học} (Curriculum) không chỉ là danh sách chủ đề phải dạy, mà là cách ý tưởng toán học được lựa chọn, sắp xếp và liên kết thành lộ trình học tập. Nó trả lời câu hỏi: học sinh sẽ được gặp ý tưởng nào, theo thứ tự nào, và với mức độ sâu ra sao.
      
      \item \emph{Giảng dạy} (Instruction) không chỉ là việc giáo viên trình bày kiến thức, mà là toàn bộ cách tổ chức việc học trong lớp: chọn nhiệm vụ nào, đặt câu hỏi ra sao, điều phối thảo luận thế nào, và tạo cơ hội nào để học sinh tham gia vào luận và hiểu.
      
      \item \emph{Đánh giá} (Assessment) không chỉ là bài kiểm tra cho điểm, mà là mọi cách nhận biết và ghi nhận việc học của học sinh. Nó cho thấy điều gì được coi là quan trọng, đồng thời tác động ngược trở lại đến cách dạy và cách học.
    \end{itemize}\par\vspace*{\baselineskip}

    \textbf{Ý nghĩa sư phạm}\par\vspace*{\baselineskip}

    Điểm cốt lõi là: không thể mong học sinh phát triển luận và hiểu nếu chương trình học, giảng dạy và đánh giá gửi đi tín hiệu khác nhau. Chẳng hạn, nếu lớp học khuyến khích học sinh giải thích và kết nối ý tưởng, nhưng đánh giá chủ yếu chỉ đòi hỏi thực hiện đúng thủ tục, thì điều được củng cố sau cùng sẽ không phải là luận và hiểu. Vì vậy, tính nhất quán ở đây không phải là phẩm chất phụ thêm vào, mà là điều kiện cấu trúc để luận và hiểu trở thành một phần thực sự của trải nghiệm học toán.\par\vspace*{\baselineskip}

    \textbf{Liên hệ với mạch sách}\par\vspace*{\baselineskip}

    Nếu Chương 9 đặt trọng tâm vào câu hỏi ai có cơ hội được tham gia vào luận và hiểu, thì Chương 10 chuyển sang câu hỏi những cơ hội ấy có được hệ thống hóa và nâng đỡ một cách nhất quán hay không. Trọng tâm vì thế được chuyển từ vấn đề công bằng trong tiếp cận sang vấn đề tổ chức chương trình học, giảng dạy và đánh giá.}
  }

\TriPara
  {
    \EN{\noindent To achieve the vision set forth in this publication, the components of the educational system must work together for the good of the students served by the system. Curriculum, instruction, and assessment can be designed to support students' reasoning and sense making. A coherent and cohesive mathematics program requires strong alignment of these three elements.}
  }
  {
    \VI{\noindent Để hiện thực hóa tầm nhìn được nêu ra trong ấn phẩm này, các thành tố của hệ thống giáo dục phải phối hợp với nhau vì lợi ích của học sinh được hệ thống ấy phục vụ. Chương trình học, giảng dạy và đánh giá có thể được thiết kế sao cho hỗ trợ luận và hiểu của học sinh. Chương trình toán học mạch lạc và gắn kết đòi hỏi sự ăn khớp chặt chẽ giữa ba thành tố này.}
  }{\NOTE{}}

\ParallelSection
  {Curriculum and Instruction}
  {Chương trình học và giảng dạy}
  {}

\TriPara
  {
    \EN{Increasingly, calls are being made for consistency in the high school mathematics curriculum across the United States (National Science Board 2007). Several sets of recommendations propose a list of content topics delineated by grade level, reflecting a range of goals and priorities (American Diploma Project 2004; College Board 2006, 2007; ACT 2007; Achieve 2007a, 2007b; Franklin et al. 2007). This publication calls for a different type of consistency across the curriculum---curricular emphases and instructional approaches that make reasoning and sense making foundational to the content that is taught and learned. Along with the more-detailed content recommendations outlined in \emph{Principles and Standards} (NCTM 2000a), this publication provides a critical filter in examining any curriculum arrangement to ensure that the ultimate goals of the high school mathematics program are achieved.}
  }
  {
    \VI{Ngày càng có nhiều lời kêu gọi về tính nhất quán trong chương trình toán trung học phổ thông trên toàn Hoa Kì (National Science Board 2007). Một số bộ khuyến nghị đề xuất danh sách chủ đề nội dung được phân định theo từng khối lớp, phản ánh nhiều mục tiêu và ưu tiên khác nhau (American Diploma Project 2004; College Board 2006, 2007; ACT 2007; Achieve 2007a, 2007b; Franklin et al. 2007). Tuy nhiên, ấn phẩm này kêu gọi kiểu nhất quán khác trong toàn bộ chương trình học---đó là trọng tâm chương trình học và cách tiếp cận giảng dạy làm cho luận và hiểu trở thành nền tảng của nội dung được dạy và được học. Cùng với những khuyến nghị nội dung chi tiết hơn được trình bày trong \emph{Principles and Standards} (NCTM 2000a), ấn phẩm này cung cấp bộ lọc mang tính phê phán để xem xét bất kì cách sắp xếp chương trình nào, nhằm bảo đảm rằng mục tiêu tối hậu của chương trình toán trung học phổ thông được hiện thực hóa.}
  }{\NOTE{}}

\TriPara
  {
    \EN{Classrooms that promote the goals of this publication prepare students to view mathematics as a connected whole across mathematical domains that consistently involves reasoning and sense making. Fostering depth of students' mathematical knowledge requires classrooms in which students are actively involved in solving problems that require them to make connections among content areas and to develop mathematical reasoning habits. The selection of mathematical tasks to use in instruction is central in promoting reasoning and sense making. Consider some of the qualities of the tasks given in the examples of this publication, including whether they---

    \begin{itemize}
      \item promote sound and significant mathematical content;
      
      \item reflect students' understandings, interests, and experiences;
      
      \item support the range of ways that diverse students learn mathematics;
      
      \item engage students' intellect by requiring reasoning and problem solving;
      
      \item help students build connections; and
      
      \item promote communication (cf. NCTM [1991]).
    \end{itemize}

    Without engaging tasks, teachers' efforts to promote students' reasoning will be limited, no matter how skilled their efforts. The approach common in the United States for many years (Martin 2007), in which teachers introduce a new topic, work out several examples, and then ask their students to emulate what they have seen, will not achieve the goal. A more profitable approach may be ``teaching via problem solving'' (Schroder and Lester 1989), in which teachers use engaging tasks to help students reason about, and make sense of, new mathematical ideas.}
  }
  {
    \VI{Lớp học hướng tới mục tiêu của ấn phẩm này giúp học sinh nhìn nhận toán học như một chỉnh thể liên kết xuyên suốt các lĩnh vực toán học, trong đó luận và hiểu luôn hiện diện một cách nhất quán. Việc bồi dưỡng chiều sâu tri thức toán học của học sinh đòi hỏi lớp học trong đó các em tham gia tích cực vào việc giải quyết những bài toán buộc các em phải tạo kết nối giữa các mảng nội dung và phát triển thói quen luận toán học. Việc lựa chọn nhiệm vụ toán học dùng trong giảng dạy giữ vai trò trung tâm trong việc thúc đẩy luận và hiểu. Hãy xét một số phẩm chất của nhiệm vụ xuất hiện trong ví dụ của ấn phẩm này, chẳng hạn liệu chúng có---\vspace*{\baselineskip}

    \begin{itemize}
      \item thúc đẩy nội dung toán học vững chắc và quan trọng;
      
      \item phản ánh sự hiểu biết, mối quan tâm và trải nghiệm của học sinh;
      
      \item hỗ trợ nhiều cách học toán khác nhau của học sinh đa dạng;
      
      \item huy động trí tuệ của học sinh bằng cách đòi hỏi luận và giải quyết vấn đề;
      
      \item giúp học sinh xây dựng kết nối; và
      
      \item thúc đẩy giao tiếp (x. NCTM [1991]).
    \end{itemize}

    Nếu không có nhiệm vụ hấp dẫn, nỗ lực của giáo viên nhằm thúc đẩy luận của học sinh sẽ bị hạn chế, cho dù giáo viên có khéo léo đến đâu trong nỗ lực ấy. Cách tiếp cận đã phổ biến tại Hoa Kì trong nhiều năm (Martin 2007)---trong đó giáo viên giới thiệu chủ đề mới, làm mẫu một số ví dụ, rồi yêu cầu học sinh bắt chước những gì đã thấy---sẽ không đạt được mục tiêu này. Cách tiếp cận hiệu quả hơn có thể là ``dạy học thông qua giải quyết vấn đề'' (Schroder và Lester 1989), trong đó giáo viên sử dụng nhiệm vụ hấp dẫn để giúp học sinh luận về và hiểu ý tưởng toán học mới.}
  }{\NOTE{}}

\TriPara
  {
    \EN{As stated in the Curriculum Principle, ``a curriculum is more than a collection of activities: it must be coherent, focused on important mathematics ...'' (NCTM 2000a, p. 14). Careful attention must be given to developing sequences of tasks that help students learn important mathematics. Integrated curricula are particularly well suited to meet this goal because they are designed to emphasize connections across mathematical domains. Regardless of the way a curriculum is organized, those involved in the mathematics program must do their best to ensure that all students are being given a strong preparation in mathematics that emphasizes the curricular connections needed to support effective mathematical reasoning and sense making.}
  }
  {
    \VI{Như đã nêu trong Nguyên lí Chương trình học, ``chương trình học không chỉ là tập hợp các hoạt động: nó phải nhất quán, tập trung vào toán học quan trọng ...'' (NCTM 2000a, trang 14). Cần chú ý kĩ lưỡng đến việc xây dựng các chuỗi nhiệm vụ giúp học sinh học toán học quan trọng. Chương trình học tích hợp đặc biệt phù hợp với mục tiêu này vì được thiết kế để nhấn mạnh mối liên hệ giữa các lĩnh vực toán học. Bất kể chương trình được tổ chức theo cách nào, những người tham gia vào chương trình toán học đều phải nỗ lực hết sức để bảo đảm rằng mọi học sinh đều được trang bị nền tảng toán học vững chắc, trong đó nhấn mạnh những kết nối trong chương trình học cần thiết để hỗ trợ luận toán học và hiểu toán học một cách hiệu quả.}
  }{\NOTE{}}

\TriPara
  {
    \EN{Merely creating an organization of worthwhile mathematical tasks within a coherent curriculum is clearly not enough. The tasks must be carried out in a classroom in which a teacher has built a classroom culture that values students' reasoning and sense making. Effective instruction requires having and communicating high expectations for all students. In such learning environments, students continually explore interesting tasks, either individually or in groups, and communicate their conjectures and conclusions to others. Pedagogical techniques and carefully designed instructional materials further ensure that the focus of instruction is on students' thinking and reasoning, and provide opportunities for all students to participate. Professional development is necessary to ensure that teachers have the tools they need to build such a classroom culture. See ``Developing Reasoning Habits in the Classroom'' in chapter 2 for additional commentary on teaching practices that may help develop reasoning and sense making.}
  }
  {
    \VI{Rõ ràng, chỉ xây dựng một hệ thống nhiệm vụ toán học có giá trị trong chương trình học nhất quán là chưa đủ. Những nhiệm vụ ấy cần được triển khai trong lớp học nơi giáo viên đã xây dựng được văn hóa coi trọng luận và hiểu của học sinh. Việc giảng dạy hiệu quả đòi hỏi giáo viên phải đặt ra và truyền đạt kì vọng cao đối với tất cả học sinh. Trong môi trường học tập như vậy, học sinh liên tục khám phá những nhiệm vụ thú vị, cá nhân hoặc theo nhóm, và trao đổi với người khác về phỏng đoán cũng như kết luận của mình. Kĩ thuật sư phạm và học liệu được thiết kế cẩn trọng còn giúp bảo đảm rằng trọng tâm giảng dạy nằm ở tư duy và luận của học sinh, đồng thời tạo cơ hội để mọi học sinh đều có thể tham gia. Hoạt động phát triển chuyên môn là cần thiết để bảo đảm rằng giáo viên có được những công cụ cần thiết để xây dựng văn hóa lớp học như vậy. Có thể xem thêm mục ``Phát triển thói quen luận trong lớp học'' trong Chương 2 để tham khảo thêm bình luận về những thực hành giảng dạy có thể giúp phát triển luận và hiểu.}
  }{\NOTE{}}

\ParallelSection
  {Vertical Alignment of Curriculum}
  {Sự liên thông theo chiều dọc của chương trình học}
  {}

\TriPara
  {
    \EN{All too often students experience discontinuities as they move from one level of mathematics to another (Smith et al. 2000), from middle school to high school and from high school to postsecondary education. The alignment of mathematical expectations related to the goals of this publication is important as students move from one school level to another. To achieve the goal of curricular coherence, an open dialogue is essential among prekindergarten through grade 8 teachers, high school mathematics teachers, mathematics teacher educators, and mathematics and statistics faculty in higher education, as well as from client disciplines, to ensure continuing support and development of students' mathematical abilities.}
  }
  {
    \VI{Học sinh thường xuyên trải qua những đứt gãy khi chuyển từ cấp học toán này sang cấp học toán khác (Smith et al. 2000), từ trung học cơ sở lên trung học phổ thông và từ trung học phổ thông lên giáo dục sau trung học. Sự liên kết giữa các kì vọng toán học gắn với mục tiêu của ấn phẩm này là hết sức quan trọng khi học sinh chuyển từ cấp học này sang cấp học khác. Để đạt được mục tiêu về tính nhất quán của chương trình học, cần có đối thoại cởi mở giữa giáo viên từ mầm non đến lớp 8, giáo viên toán trung học phổ thông, nhà đào tạo giáo viên toán, giảng viên toán và thống kê trong giáo dục đại học, cũng như giảng viên từ các ngành học sử dụng toán, nhằm bảo đảm sự hỗ trợ và phát triển liên tục năng lực toán học của học sinh.}
  }{\NOTE{}}

\TriPara
  {
    \EN{Schools, families, policymakers, and others need to see evidence that reasoning and sense making abilities are shared goals among all levels of mathematics teaching, including elementary school mathematics, high school mathematics, and the undergraduate curriculum. The time is right to build strong partnerships among the various constituencies, and to recognize the potential beneficial consequences resulting from a mathematics curriculum that focuses on reasoning and sense making as articulated parts of the continuum from prekindergarten through grade 16.}
  }
  {
    \VI{Nhà trường, gia đình, nhà hoạch định chính sách và các bên liên quan khác cần nhìn thấy bằng chứng rõ ràng rằng năng lực luận và hiểu là mục tiêu chung được chia sẻ ở mọi cấp độ dạy toán, bao gồm toán tiểu học, toán trung học phổ thông và chương trình đại học. Đây là thời điểm thích hợp để xây dựng quan hệ hợp tác vững chắc giữa các nhóm liên quan khác nhau, đồng thời nhận ra những hệ quả tích cực tiềm tàng của chương trình toán học đặt luận và hiểu làm trọng tâm, với hai yếu tố này được nêu rõ như những thành phần của trục liên tục từ tiền mẫu giáo đến lớp 16.}
  }
  {
    \NOTE{Cụm ``từ tiền mẫu giáo đến lớp 16'' phản ánh cách nói trong hệ giáo dục Hoa Kì, trong đó lớp 16 thường tương ứng với năm cuối bậc đại học bốn năm. Ở đây, điểm nhấn không nằm ở cách chia cấp lớp, mà ở yêu cầu liên kết mục tiêu luận và hiểu trên toàn bộ trục học tập, từ trước tiểu học đến giáo dục đại học.}
  }

\TriPara
  {
    \EN{
      \textbf{\emph{Alignment with prekindergarten through grade 8 mathematics}}\vspace{\baselineskip}

      This publication takes the recommendations of Curriculum Focal Points for Prekindergarten through Grade 8 Mathematics (Curriculum Focal Points) (NCTM 2006a) as its starting point for ensuring that proper alignment occurs between the middle grades and high school. For each grade level, prekindergarten through grade 8, Curriculum Focal Points outlines Focal Point content areas and describes objectives in additional areas of ``related connections.'' Each chapter of Curriculum Focal Points begins with two reminders:\vspace{\baselineskip}

      \begin{itemize}
        \item The ``curriculum focal points and related connections ... are the recommended content emphases for mathematics....'' That is, the objectives outlined in the additional ``related connections'' are also required material. Otherwise, students will not receive the breadth of mathematical content required for high school, particularly in the areas of data analysis, probability, and statistics.
        
        \item ``It is essential that these focal points be addressed in contexts that promote problem solving, reasoning, communication, making connections, and designing and analyzing representations.'' That is, teaching skills without providing a foundation in, and without giving attention to, important mathematical processes, including reasoning and sense making, will not adequately prepare students for the kinds of thinking expected in high school mathematics.
      \end{itemize}

      All students entering high school should have a sound base of mathematical experiences as described in the broadest interpretation of \emph{Curriculum Focal Points}. 

      Taking algebra in middle grades has become increasingly popular (National Center for Educational Statistics 2009). This publication contends that whenever new content is encountered-in elementary, middle, or high school-the foundation for conceptual understanding is laid, in part, through opportunities for students to reason with and about the mathematics being studied and to connect it with their prior knowledge. Therefore, middle school courses that contain significant algebraic content must build on and foster students' reasoning and sense making-including making connections across various mathematical domains-to lay a foundation for success in later courses. For students without the content background described in Curriculum Focal Points and with insufficient experience in reasoning and sense making, taking a formal algebra course prematurely will result in missed opportunities for broadening and strengthening their understanding of fundamental mathematical concepts, including those that are foundational for later success in algebra. Taking a formal algebra course too early will be counterproductive unless students have the necessary prerequisite skills because ``exposing students to such coursework before they are ready often leads to frustration, failure, and negative attitudes toward mathematics and learning'' (NCTM 2008).}
  }
  {
    \VI{
      \textbf{\emph{Sự liên kết với toán học từ tiền mẫu giáo đến lớp 8}}\vspace{\baselineskip}

      Ấn phẩm này lấy khuyến nghị của \emph{Curriculum Focal Points for Prekindergarten through Grade 8 Mathematics (Curriculum Focal Points)} (NCTM 2006a) làm điểm xuất phát để bảo đảm sự liên kết thích đáng giữa các lớp trung học cơ sở và trung học phổ thông. Với mỗi cấp lớp từ tiền mẫu giáo đến lớp 8, \emph{Curriculum Focal Points} nêu ra các mảng nội dung trọng tâm và mô tả mục tiêu trong các mảng bổ sung thuộc phần ``kết nối liên quan''. Mỗi chương của \emph{Curriculum Focal Points} đều mở đầu bằng hai lời nhắc:\vspace{\baselineskip}

    \begin{itemize}
      \item ``Các trọng tâm chương trình học và kết nối liên quan ... là những trọng tâm nội dung được khuyến nghị cho môn toán....'' Nói cách khác, mục tiêu được nêu trong phần bổ sung ``kết nối liên quan'' cũng là nội dung bắt buộc. Nếu không, học sinh sẽ không có được bề rộng nội dung toán học cần thiết cho trung học phổ thông, đặc biệt trong các mảng phân tích dữ liệu, xác suất và thống kê.
      
      \item ``Điều thiết yếu là những trọng tâm này phải được xử lí trong các bối cảnh thúc đẩy giải quyết vấn đề, luận, giao tiếp, thiết lập kết nối, và thiết kế cũng như phân tích biểu diễn.'' Nói cách khác, nếu chỉ dạy kĩ năng mà không tạo nền tảng cho, cũng không chú ý đến, những quá trình toán học quan trọng, trong đó có luận và hiểu, thì học sinh sẽ không được chuẩn bị đầy đủ cho kiểu suy nghĩ được mong đợi trong toán học trung học phổ thông.
    \end{itemize}\vspace*{\baselineskip}

    Mọi học sinh bước vào trung học phổ thông đều cần có nền tảng trải nghiệm toán học vững chắc theo cách hiểu rộng nhất của \emph{Curriculum Focal Points}.

    Việc học đại số ở bậc trung học cơ sở ngày càng trở nên phổ biến (National Center for Educational Statistics 2009). Ấn phẩm này cho rằng, mỗi khi học sinh gặp nội dung mới---ở tiểu học, trung học cơ sở hay trung học phổ thông---thì nền tảng cho hiểu khái niệm được hình thành, một phần, nhờ cơ hội để các em vừa dùng toán học đang học để luận, vừa luận về chính toán học ấy, đồng thời kết nối nó với kiến thức đã có. Vì vậy, khóa học ở bậc trung học cơ sở có nội dung đại số đáng kể phải dựa trên và đồng thời nuôi dưỡng luận và hiểu của học sinh---bao gồm cả việc thiết lập các kết nối giữa các lĩnh vực toán học khác nhau---để đặt nền tảng cho thành công trong các khóa học về sau. Đối với học sinh chưa có nền tảng nội dung như được mô tả trong \emph{Curriculum Focal Points} và cũng chưa có đủ trải nghiệm về luận và hiểu, việc học khóa học đại số chính thức quá sớm sẽ khiến các em bỏ lỡ cơ hội để mở rộng và củng cố sự hiểu của mình về các khái niệm toán học cơ bản, trong đó có những khái niệm nền tảng cho thành công sau này trong đại số. Việc học khóa học đại số chính thức quá sớm sẽ phản tác dụng nếu học sinh chưa có kĩ năng tiên quyết cần thiết, bởi ``việc cho học sinh tiếp xúc với loại khóa học như vậy trước khi các em sẵn sàng thường dẫn đến nản lòng, thất bại, và thái độ tiêu cực đối với toán học và việc học'' (NCTM 2008).}
  }{\NOTE{}}

\TriPara
  {
    \EN{
      \textbf{\emph{Alignment with postsecondary education}}\vspace{\baselineskip}

      \emph{Principles and Standards} (NCTM 2000a) states, ``All students are expected to study mathematics each of the four years that they are enrolled in high school'' (p. 288), and students must experience reasoning and sense making across all four years of a high-quality high school mathematics program. This expectation is reinforced in NCTM's (2006b) Position Statement on Time: ``Evidence supports the enrollment of high school students in a mathematics course every year, continuing beyond the equivalent of a second year of algebra and a year of geometry.''

      However, what courses to offer students who have taken formal algebra (and perhaps even geometry) in middle school has become increasingly problematic. If students take calculus in high school, they should take rigorous courses following the Advanced Placement guidelines (College Board 2008a) so that they are prepared to continue their study of mathematics at the collegiate level. Advanced Placement Statistics has become increasingly popular (College Board 2008b) and will meet the needs of many students going on to either mathematics-related or non-mathematics-related majors in college. Developing alternative courses that meet the needs of other students should also be considered---for example, capstone courses are being developed to help students connect what they have been learning with their future goals (cf. Charles A. Dana Center 2008). All upper-level courses should continue to promote the goals of reasoning and sense making as set forth in this publication.

      The recommendations in this publication are consistent with recommendations by the Mathematics Association of America (Ganter and Barker 2004, the Committee on the Undergraduate Program in Mathematics [CUPM] 2004), and the American Mathematical Association of Two-Year Colleges (2006). These organizations propose that the entire college-level mathematics curriculum, for all students, even those who take just one course, should develop ``analytical, critical reasoning, problem-solving, and communication skills...'' (CUPM 2004, pp. 5, 13). The commonalities between recommendations for the undergraduate curriculum and those of the current publication provide opportunities to strengthen articulation between high school and postsecondary education. Furthermore, the framers of this publication believe that focusing on reasoning and sense making can move us forward on important issues, such as those involving mathematics placement examinations used by postsecondary institutions and articulation on the role of technology in high school and college mathematics.}
  }
  {
    \VI{
    \textbf{\emph{Sự liên kết với giáo dục sau trung học}}\vspace{\baselineskip}

      \emph{Principles and Standards} (NCTM 2000a) nêu rõ: ``Tất cả học sinh được kì vọng sẽ học toán trong cả bốn năm các em theo học trung học phổ thông'' (trang 288), và học sinh cần trải nghiệm luận và hiểu xuyên suốt cả bốn năm của chương trình toán trung học phổ thông chất lượng cao. Kì vọng này được củng cố trong Tuyên bố lập trường về Thời lượng của NCTM (2006b): ``Bằng chứng ủng hộ việc học sinh trung học phổ thông học một khóa toán mỗi năm, tiếp tục sau mức tương đương với năm thứ hai đại số và một năm hình học.''

      Tuy nhiên, việc nên mở những khóa học nào cho học sinh đã học đại số chính thức (và có thể cả hình học) ở trung học cơ sở ngày càng trở nên nan giải. Nếu học sinh học giải tích ở trung học phổ thông, các em cần theo học các khóa học nghiêm ngặt theo hướng dẫn của chương trình Advanced Placement (College Board 2008a) để được chuẩn bị đầy đủ cho việc tiếp tục học toán ở bậc đại học. Advanced Placement Statistics ngày càng trở nên phổ biến (College Board 2008b) và đáp ứng nhu cầu của nhiều học sinh dự định theo học ngành có liên quan hoặc không liên quan trực tiếp đến toán học ở bậc đại học. Việc phát triển các khóa học thay thế nhằm đáp ứng nhu cầu của những học sinh khác cũng cần được xem xét---chẳng hạn, các khóa học tổng kết đang được xây dựng để giúp học sinh kết nối những gì các em đã học với mục tiêu tương lai (xem Charles A. Dana Center 2008). Tất cả khóa học nâng cao đều cần tiếp tục thúc đẩy mục tiêu luận và hiểu như đã nêu trong ấn phẩm này.

      Các khuyến nghị trong ấn phẩm này phù hợp với khuyến nghị của Mathematics Association of America (Ganter and Barker 2004; Ủy ban Chương trình Toán học bậc Đại học [CUPM] 2004) và American Mathematical Association of Two-Year Colleges (2006). Các tổ chức này đề xuất rằng toàn bộ chương trình toán học bậc đại học, dành cho tất cả sinh viên, kể cả người chỉ học một môn toán duy nhất, đều cần phát triển ``kĩ năng phân tích, luận phản biện, giải quyết vấn đề và giao tiếp...'' (CUPM 2004, trang 5, 13). Những điểm tương đồng giữa khuyến nghị cho chương trình đại học và khuyến nghị của ấn phẩm hiện tại mở ra cơ hội tăng cường sự liên kết giữa trung học phổ thông và giáo dục sau trung học. Hơn nữa, những người xây dựng ấn phẩm này tin rằng việc tập trung vào luận và hiểu có thể giúp thúc đẩy tiến bộ trong các vấn đề quan trọng, chẳng hạn như kì thi xếp lớp toán được các cơ sở giáo dục sau trung học sử dụng và việc thống nhất quan điểm về vai trò của công nghệ trong toán học ở trung học phổ thông và đại học.}
  }{\NOTE{}}

\ParallelSection
  {Assessment}
  {Đánh giá}
  {}

\TriPara
  {
    \EN{What we assess should reflect what we value. Assessments that support the goals of this publication will attend to students' capabilities in mathematical reasoning and sense making and contribute to students' progress in mathematics. This endeavor is essential for at least two reasons. First, we will not be able to gauge whether we are meeting our goals if those goals are not assessed. Second, high-stakes testing that concentrates primarily on procedural skills without assessing reasoning and sense making sends a message that is contrary to the vision set forth in this publication and can adversely influence instruction and students' learning. Assessment that focuses primarily on students' ability to do algebraic manipulations, apply geometric formulas, and perform basic statistical computations will lead students to believe that reasoning and sense making are not important.}
  }
  {
    \VI{Những gì chúng ta đánh giá phải phản ánh những gì chúng ta coi trọng. Các hình thức đánh giá hỗ trợ mục tiêu của ấn phẩm này cần chú ý đến năng lực luận toán học và hiểu của học sinh, đồng thời góp phần thúc đẩy sự tiến bộ của các em trong toán học. Nỗ lực này là thiết yếu vì ít nhất hai lí do. Thứ nhất, chúng ta sẽ không thể biết liệu mình có đạt được mục tiêu đã đặt ra hay không nếu mục tiêu đó không được đánh giá. Thứ hai, kì kiểm tra mang tính quyết định cao, nếu chủ yếu tập trung vào kĩ năng thủ tục mà không đánh giá luận và hiểu, sẽ gửi đi thông điệp trái ngược với tầm nhìn của ấn phẩm này, đồng thời có thể tác động tiêu cực đến giảng dạy và việc học của học sinh. Hình thức đánh giá chủ yếu nhắm vào khả năng thực hiện biến đổi đại số, áp dụng công thức hình học và thực hiện phép tính thống kê cơ bản sẽ khiến học sinh tin rằng luận và hiểu không quan trọng.}
  }{\NOTE{}}

\TriPara
  {
    \EN{Assessing students' progress in mathematics serves two distinct roles. One essential role of assessment is to provide summative measures of students' understanding of and ability to use mathematics. All too often, high-stakes tests assess minimal standards that do not reflect the content that students need for their future success. As students, educators, and the community focus on students' success in learning lower-level content, the importance of pushing students to succeed with higher-level content may be diminished. Moreover, this publication calls for further progress in the design of many high-stakes test items to better assess students' reasoning and sense making.}
  }
  {
    \VI{Việc đánh giá sự tiến bộ của học sinh trong toán học có hai vai trò khác biệt. Một vai trò thiết yếu của đánh giá là cung cấp các thước đo tổng kết về sự hiểu toán học và khả năng sử dụng toán học của học sinh. Rất thường xuyên, các kì kiểm tra mang tính quyết định cao chỉ đánh giá chuẩn tối thiểu, vốn không phản ánh nội dung học sinh cần cho thành công trong tương lai. Khi học sinh, nhà giáo dục và cộng đồng cùng tập trung vào thành công của học sinh trong việc học nội dung ở mức thấp hơn, tầm quan trọng của việc thúc đẩy học sinh thành công với nội dung ở mức cao hơn có thể bị xem nhẹ. Hơn nữa, ấn phẩm này kêu gọi tiếp tục cải thiện thiết kế nhiều câu hỏi trong các kì kiểm tra mang tính quyết định cao, nhằm đánh giá tốt hơn năng lực luận và hiểu của học sinh.}
  }{\NOTE{}}

\TriPara
  {
    \EN{The second essential component of assessment is formative assessment, which is an integral part of the learning process for the student. According to Black and Wiliam (1998), formative assessment involves providing students with learning activities and, on the basis of feedback from those activities, adjusting teaching to meet the students' needs. Various forms of formative assessment can furnish information about students' reasoning and sense making. These include teacher observations, classroom discussions, student journals, student presentations, homework, and inclass tasks, as well as tests and quizzes that ask students to explain their thinking. To ensure that their assessments reflect curricular priorities, teachers must use both summative and formative instruments for assessing students' understanding of mathematics. Example 3, ``Around Pi,'' illustrates the use of an in-class assignment for formative assessment. Regardless of whether a grade is assigned, the teacher might design a rubric, such as that in figure 10.1, to better understand the students' progress in understanding error and relative error.}
  }
  {
    \VI{Thành phần thiết yếu thứ hai của đánh giá là đánh giá quá trình, vốn là phần không thể tách rời của quá trình học tập của học sinh. Theo Black and Wiliam (1998), đánh giá quá trình bao gồm việc cung cấp cho học sinh hoạt động học tập và, dựa trên phản hồi từ các hoạt động đó, điều chỉnh giảng dạy để đáp ứng nhu cầu của học sinh. Nhiều hình thức đánh giá quá trình có thể cung cấp thông tin về luận và hiểu của học sinh. Các hình thức này bao gồm quan sát của giáo viên, thảo luận trên lớp, nhật kí học tập của học sinh, bài trình bày của học sinh, bài tập về nhà, nhiệm vụ trên lớp, cũng như bài kiểm tra và bài kiểm tra ngắn yêu cầu học sinh giải thích suy nghĩ của mình. Để bảo đảm việc đánh giá phản ánh các ưu tiên của chương trình học, giáo viên phải dùng cả công cụ đánh giá tổng kết lẫn đánh giá quá trình để đánh giá sự hiểu toán học của học sinh. Ví dụ 3, ``Số Pi dài - ngắn thế nào?'', minh họa việc dùng nhiệm vụ trên lớp cho đánh giá quá trình. Dù có cho điểm hay không, giáo viên có thể thiết kế bảng tiêu chí đánh giá, chẳng hạn như bảng trong Hình 10.1, để hiểu rõ hơn tiến triển của học sinh trong sự hiểu về sai số và sai số tương đối.}
  }
  {
    \NOTE{Đánh giá quá trình không chỉ nhằm ghi nhận kết quả học tập, mà còn nhằm tạo phản hồi để giáo viên điều chỉnh dạy học và học sinh tiếp tục phát triển hiểu biết.}
  }

\TriFigure
  {../figures/fhsm_2009_p105_f1.pdf}
  {../figures/fhsm_2009_p105_f1_vi.pdf}
  {A rubric for assessing responses to example 3}
  {Bảng tiêu chí đánh giá câu trả lời cho Ví dụ 3.}
  {fig:fig_10.1}

\ParallelSection
  {Conclusion}
  {Kết luận}
  {}

\TriPara
  {
    \EN{Alignment of curriculum, instruction, and assessment is a starting point on the road to implementing the goals set forth in this publication within the large and complex educational system of the United States. These three components cannot be viewed in isolation. Rather, decision makers or policymakers dealing with aspects of any one of these areas must ask themselves, ``Is a coherent mathematics program being developed that takes into account all three programmatic elements?''}
  }
  {
    \VI{Sự liên kết giữa chương trình học, giảng dạy và đánh giá là điểm khởi đầu trên con đường hiện thực hóa mục tiêu được nêu ra trong ấn phẩm này trong hệ thống giáo dục lớn và phức tạp của Hoa Kì. Không thể xem ba thành tố này một cách tách rời. Trái lại, người ra quyết định hoặc nhà hoạch định chính sách xử lí bất kì phương diện nào của một trong ba thành tố ấy đều phải tự hỏi: ``Liệu có đang xây dựng chương trình toán học nhất quán, có tính đến cả ba thành tố chương trình hay không?''}
  }{\NOTE{}}

\TriPara
  {
    \EN{Efforts to focus the mathematics curriculum on reasoning and sense making will require responsible efforts of, and considerable work from, all stakeholders. Support will be required through governing policy statements that clearly align with the goals set forth here. The realization of these recommendations will require considerable time and resources to provide the significant professional development required for school faculty and administrators. Many of these efforts are discussed further in the next chapter.}
  }
  {
    \VI{Nỗ lực nhằm đặt luận và hiểu vào trọng tâm chương trình học toán sẽ đòi hỏi nỗ lực có trách nhiệm và khối lượng công việc đáng kể từ tất cả các bên liên quan. Sự hỗ trợ sẽ cần được thể hiện qua những tuyên bố chính sách quản lí phù hợp rõ ràng với các mục tiêu được nêu trong ấn phẩm này. Việc hiện thực hóa những khuyến nghị này sẽ cần nhiều thời gian và nguồn lực để cung cấp hoạt động phát triển chuyên môn đáng kể cho đội ngũ giáo viên và cán bộ quản lí nhà trường. Nhiều khía cạnh của nỗ lực này sẽ được bàn luận chi tiết hơn trong chương tiếp theo.}
  }{\NOTE{}}

\CloseGrid